% Biên dịch bằng XeLaTeX (2 lần để mục lục đúng):  xelatex slide_demo.tex
% Có thư mục fonts/ và assets/ cạnh file này thì dùng font + logo của template,
% không có thì tự dùng font mặc định, vẫn biên dịch được.
\documentclass[aspectratio=169,9pt]{beamer}
\usepackage{fontspec}
\usepackage[vietnamese]{babel}
\usepackage{amsmath,amssymb}
\usepackage{tikz}
\usetikzlibrary{calc,arrows.meta,positioning}
\usepackage{booktabs}
\usepackage{tabularx}
\usepackage{graphicx}
\graphicspath{{images_part3/}}
\newcolumntype{Y}{>{\raggedright\arraybackslash}X} % cột X căn trái, không giãn chữ

% ── Fonts: Bungee (tiêu đề) + Be Vietnam Pro (nội dung), có fallback ──
\IfFileExists{fonts/BeVietnamPro-Regular.ttf}{%
  \setmainfont{BeVietnamPro-Regular.ttf}[
    Path=fonts/,
    BoldFont=BeVietnamPro-Bold.ttf,
    ItalicFont=BeVietnamPro-Medium.ttf,
    BoldItalicFont=BeVietnamPro-SemiBold.ttf]
  \newfontfamily\bungee{Bungee-Regular.ttf}[Path=fonts/]
  \newfontfamily\bevnsemi{BeVietnamPro-SemiBold.ttf}[Path=fonts/]
}{%
  \newcommand{\bungee}{\bfseries}
  \newcommand{\bevnsemi}{\bfseries}
}

% Logo chỉ chèn khi file tồn tại
\newcommand{\logoimg}[1]{\IfFileExists{assets/hoids1_logo.png}{\includegraphics[height=#1]{assets/hoids1_logo.png}}{}}

% ── Chèn ảnh có KHUNG GIỮ CHỖ ────────────────────────────────
% \slideimg[chiều rộng]{tên-file.png}{mô tả ảnh cần chụp}
% File ảnh để trong thư mục images/. Chưa có ảnh thì hiện khung xám ghi sẵn tên file cần bỏ vào,
% có ảnh rồi thì tự hiện ảnh (không cần sửa gì thêm).
\newcommand{\slideimg}[3][0.9\textwidth]{%
  \IfFileExists{images/#2}{%
    \includegraphics[width=#1,height=0.5\textheight,keepaspectratio]{images/#2}%
  }{%
    \fcolorbox{skyblue}{white}{\parbox[c][0.28\textheight][c]{#1}{\centering\scriptsize\color{inkblue}%
      \textbf{[CHƯA CÓ ẢNH]}\\[2pt]\texttt{images/#2}\\[2pt]#3}}%
  }%
}

% ── Bảng màu ──
\definecolor{inkblue}{HTML}{1B2C8C}
\definecolor{brandblue}{HTML}{0A6FBA}
\definecolor{midblue}{HTML}{2E7DD1}
\definecolor{skyblue}{HTML}{73A5E1}
\definecolor{paleblue}{HTML}{D4E6FB}
\definecolor{bgblue}{HTML}{EAF2FD}
\definecolor{gold}{HTML}{FFC21E}
\definecolor{footernavy}{HTML}{14276B}

\usetheme{default}
\usefonttheme{professionalfonts}
\setbeamertemplate{navigation symbols}{}
\setbeamercolor{normal text}{fg=footernavy}
\setbeamercolor{structure}{fg=brandblue}
\setbeamercolor{itemize item}{fg=brandblue}
\setbeamercolor{itemize subitem}{fg=skyblue}
\setbeamertemplate{itemize item}{\small\textbullet}

\setbeamercolor{block title}{bg=inkblue,fg=white}
\setbeamercolor{block body}{bg=white,fg=footernavy}
\setbeamercolor{block title alerted}{bg=brandblue,fg=white}
\setbeamercolor{block body alerted}{bg=white,fg=footernavy}
\setbeamercolor{block title example}{bg=gold,fg=footernavy}
\setbeamercolor{block body example}{bg=white,fg=footernavy}
\setbeamerfont{block title}{family=\bevnsemi,size=\normalsize}
\setbeamerfont{block title example}{family=\bevnsemi,size=\normalsize}
\setbeamerfont{block title alerted}{family=\bevnsemi,size=\normalsize}
\setbeamertemplate{blocks}[rounded][shadow=true]

% ── Nền mọi slide ──
\setbeamertemplate{background}{%
  \begin{tikzpicture}
    \useasboundingbox (0,0) rectangle (\the\paperwidth,\the\paperheight);
    \fill[bgblue] (0,0) rectangle (\the\paperwidth,\the\paperheight);
    \foreach \x/\h in {0/0.9, 0.9/1.5, 1.7/1.0, 2.5/1.9, 3.5/1.2, 4.3/1.7,
                       5.2/1.0, 6.0/1.6, 6.9/1.1, 7.7/2.0, 8.7/1.3, 9.6/1.7,
                       10.5/1.0, 11.3/1.5, 12.2/1.1, 13.0/1.8, 14.0/1.2, 14.9/1.6}{
      \fill[paleblue] (\x,0) rectangle (\x+0.78,\h*0.5);
    }
    \fill[footernavy] (0,0) rectangle (\the\paperwidth,0.6);
    \fill[gold] (0,0.6) rectangle (\the\paperwidth,0.65);
  \end{tikzpicture}%
}

% ── Tiêu đề frame ──
\setbeamertemplate{frametitle}{%
  \nointerlineskip%
  \vskip -2pt%
  \begin{beamercolorbox}[wd=\paperwidth,ht=1.1cm,dp=0.2cm]{frametitle}
  \hbox to \paperwidth{%
    \hspace{0.6em}%
    \begin{minipage}[c]{0.80\paperwidth}
      {\bungee\color{inkblue}\fontsize{11}{11}\selectfont \insertframetitle}%
      \par\nointerlineskip\vspace{2pt}%
      {\color{brandblue}\rule{\linewidth}{1.2pt}}
    \end{minipage}%
    \hfill%
    \raisebox{-4pt}{\logoimg{0.085\paperheight}}%
    \hspace{0.6em}%
  }%
  \end{beamercolorbox}%
  \vspace{-0.6em}%
}

% ── Chân trang ──
\setbeamertemplate{footline}{%
  \begin{beamercolorbox}[wd=\paperwidth,ht=0.6cm,dp=0pt]{footline}
    \vbox to 0.6cm{%
      \vfil
      \hbox to \paperwidth{%
        \hspace{1.2em}%
        {\footnotesize\color{white}\bevnsemi NT106 $\cdot$ Battleship Network}\hfill
        {\footnotesize\color{white}\insertframenumber{} / \inserttotalframenumber}\hspace{1.2em}%
      }%
      \vfil
    }%
  \end{beamercolorbox}%
}

\title{BATTLESHIP MULTIPLAYER\\KẾT NỐI SUPABASE VÀ WEB SERVICE}
\subtitle{Raw Socket TCP $\cdot$ AES-128 $\cdot$ PostgreSQL trên Cloud $\cdot$ REST API}
\author{Hoàng $\cdot$ Đức $\cdot$ Dũng}
\date{\today}

\begin{document}

% ── Slide 1: Trang tiêu đề ──
{\setbeamertemplate{footline}{}
\begin{frame}[plain]
  \centering
  \vspace{0.3cm}
  \logoimg{0.22\paperheight}\par
  \vspace{0.3cm}
  {\bungee\color{inkblue}\fontsize{17}{21}\selectfont BATTLESHIP MULTIPLAYER}\par
  \vspace{0.15cm}
  {\bungee\color{brandblue}\fontsize{13}{16}\selectfont KẾT NỐI SUPABASE VÀ WEB SERVICE}\par
  \vspace{0.25cm}
  {\bevnsemi\color{brandblue}\small Môn Lập trình mạng căn bản (NT106) $\cdot$ UIT}\par
  \vspace{0.15cm}
  {\bevnsemi\color{footernavy}\footnotesize Hoàng (thuyết trình chính) $\cdot$ Đức $\cdot$ Dũng}
\end{frame}}

% ── Slide 2: Mục lục ──
\begin{frame}{Mục Lục Nội Dung}
  \tableofcontents
\end{frame}

% =====================================================
% GHÉP CÁC PHẦN. Mỗi người chỉ sửa file phần của mình trong thư mục parts/
%   01_tongquan.tex     Hoàng
%   02_server_duc.tex   ĐỨC
%   03_csdl_dung.tex    DŨNG
%   04_cloud_hoang.tex  Hoàng
%   05_demo.tex         Hoàng
%   06_nhom.tex         Hoàng
% =====================================================
% PHẦN CỦA HOÀNG: Tổng quan đồ án. (Đức và Dũng KHÔNG sửa file này)
% =====================================================
\section{Tổng Quan Đồ Án}

\begin{frame}{Đồ Án Battleship Là Gì?}
  \begin{block}{Mục tiêu và chức năng (một phần đang phát triển)}
    \begin{itemize}
      \item Đăng ký, đăng nhập, danh sách \textbf{Online/Offline}, chat sảnh.
      \item \textbf{Tạo phòng} bằng ID, ghép trận tự động (\textbf{Matchmaking}).
      \item Xếp tàu trên lưới $10\times10$, bắn theo lượt, \textbf{Server là bên quyết định} Hit/Miss.
      \item Lưu lịch sử đấu, thắng/thua và điểm xếp hạng.
    \end{itemize}
  \end{block}
  \begin{exampleblock}{Công nghệ cốt lõi của môn NT106}
    C\# .NET 8 $\cdot$ \texttt{TcpListener/TcpClient} $\cdot$ async/await $\cdot$ framing 4-byte $\cdot$ JSON $\cdot$ AES-128 $\cdot$ BCrypt $\cdot$ EF Core
  \end{exampleblock}
\end{frame}

\begin{frame}{Tình Trạng Trước Khi Thêm Cloud}
  \begin{columns}[T]
    \begin{column}{0.52\textwidth}
      \begin{block}{Đã có}
        \begin{itemize}
          \item TCP Server bất đồng bộ + \textbf{PacketRouter} đăng ký handler.
          \item Giao thức: framing 4-byte + \textbf{Handshake} kiểm tra phiên bản.
          \item \textbf{Mã hóa AES-128} toàn bộ gói tin.
          \item Đăng nhập/đăng ký với mật khẩu băm \textbf{BCrypt}.
          \item Dữ liệu lưu \textbf{SQLite} cục bộ qua EF Core.
        \end{itemize}
      \end{block}
    \end{column}
    \begin{column}{0.44\textwidth}
      \begin{alertblock}{Hạn chế cần giải quyết}
        \begin{itemize}
          \item Dữ liệu nằm trên \textbf{một máy}, chưa dùng chung.
          \item Chưa có \textbf{web service} để truy vấn dữ liệu từ bên ngoài.
        \end{itemize}
      \end{alertblock}
    \end{column}
  \end{columns}
\end{frame}

% ============================================================
% PHẦN: KIẾN TRÚC SERVER VÀ GAME  (4 slide)
% Nội dung lấy từ file "Kiến Trúc Server & Hướng Dẫn Của Đức.pptx" (Thành viên A: Đức),
% đã ĐỐI CHIẾU với code thật trên main (TcpServer.cs, ClientSession.cs, RoomManager.cs, RoomHandler.cs).
% Chỉ ghi những gì đã hiện thực; phần chưa xong ghi rõ "đang phát triển".
% ============================================================
\section{Kiến Trúc Server Và Game}

% ---------- Slide: lõi TCP ----------
\begin{frame}{Lõi TCP Và Quản Lý Phiên}
  \begin{columns}[T]
    \begin{column}{0.52\textwidth}
      \begin{block}{Cách server xử lý kết nối}
        \begin{itemize}
          \item \texttt{TcpListener} lắng nghe bất đồng bộ, \textbf{mỗi client một Task} riêng.
          \item \textbf{Framing 4-byte}: đọc đủ \texttt{PayloadLength} rồi mới giải mã gói.
          \item \textbf{Khóa luồng gửi} \texttt{SemaphoreSlim(1,1)}: chặn hai Task ghi xen kẽ byte vào cùng một stream.
          \item \texttt{PacketRouter} chuyển gói tới từng \texttt{IPacketHandler} đã đăng ký.
        \end{itemize}
      \end{block}
    \end{column}
    \begin{column}{0.44\textwidth}
      \centering
      \begin{tikzpicture}[
        flow/.style={draw=inkblue, fill=white, rounded corners=3pt, thick, align=center,
                     minimum height=0.85cm, minimum width=5.2cm, font=\scriptsize},
        arr/.style={-{Stealth[length=2mm]}, thick, color=brandblue}]
        \node[flow, fill=paleblue] (a) at (0,0)     {1. \texttt{NetworkStream} $\rightarrow$ \texttt{PacketFraming} (4-byte)};
        \node[flow] (b) at (0,-1.45) {2. \texttt{ClientSession.OnPacketReceived}};
        \node[flow] (c) at (0,-2.9)  {3. \texttt{PacketRouter.RouteAsync(session, pkt)}};
        \node[flow, fill=paleblue] (d) at (0,-4.35) {4. \texttt{RoomHandler.HandleAsync(session, pkt)}};
        \draw[arr] (a) -- (b);
        \draw[arr] (b) -- (c);
        \draw[arr] (c) -- (d);
      \end{tikzpicture}

      \vspace{0.1cm}
      \scriptsize Luồng đi của một gói tin tại server.
    \end{column}
  \end{columns}
\end{frame}

% ---------- Slide: quản lý phòng ----------
\begin{frame}{Quản Lý Phòng Đấu}
  \begin{columns}[T]
    \begin{column}{0.50\textwidth}
      \begin{block}{\texttt{RoomManager}}
        \begin{itemize}
          \item \texttt{ConcurrentDictionary}: lưu danh sách phòng an toàn đa luồng.
          \item \textbf{Khóa riêng từng phòng}: hai người vào cùng một lúc không bị lỗi.
          \item \textbf{Host thoát} thì phòng bị giải tán và người còn lại được báo (\texttt{RoomUpdate}).
        \end{itemize}
      \end{block}
    \end{column}
    \begin{column}{0.46\textwidth}
      \footnotesize
      \renewcommand{\arraystretch}{1.25}
      \begin{tabularx}{\textwidth}{@{}Y Y@{}}
        \toprule
        \textbf{Gói tin} & \textbf{Chức năng} \\
        \midrule
        \texttt{GetRoomListRequest} & Danh sách phòng \\
        \texttt{CreateRoomRequest} & Tạo phòng bằng ID \\
        \texttt{JoinRoomRequest} & Vào phòng \\
        \texttt{LeaveRoom} & Rời phòng \\
        \bottomrule
      \end{tabularx}
      \vspace{0.15cm}

      \scriptsize \textbf{Đang phát triển:} nối chức năng Kick (đã có ở \texttt{RoomManager}) vào gói tin.
    \end{column}
  \end{columns}
\end{frame}

% ---------- Slide: state machine ----------
\begin{frame}{Máy Trạng Thái Của Phòng Đấu}
  \centering
  \vspace{0.6cm}
  \slideimg[0.97\textwidth]{12-state-machine.png}{Sơ đồ State Machine}

  \vspace{0.4cm}
  \footnotesize Hai trạng thái đầu \textbf{đã hiện thực}; ba trạng thái sau đang phát triển.
\end{frame}

% ---------- Slide: bảo mật + ngắt kết nối ----------
\begin{frame}{Bảo Mật Và Xử Lý Ngắt Kết Nối}
  \begin{columns}[T]
    \begin{column}{0.49\textwidth}
      \begin{block}{Dọn dẹp khi đứt mạng}
        \begin{itemize}
          \item Bắt \texttt{IOException} và \texttt{SocketException} \textbf{riêng biệt}.
          \item Khối \texttt{finally} luôn gọi \texttt{Cleanup()} để đóng luồng mạng.
          \item Sự kiện \texttt{OnDisconnected}: xóa khỏi danh sách online và \textbf{báo Offline} cho mọi người.
        \end{itemize}
      \end{block}
    \end{column}
    \begin{column}{0.47\textwidth}
      \begin{exampleblock}{Mã hóa đường truyền}
        \begin{itemize}
          \item Toàn bộ payload JSON qua TCP được mã hóa \textbf{AES-128}.
          \item Gói chữ rõ không có khóa bị server \textbf{đóng kết nối}.
        \end{itemize}
      \end{exampleblock}
      \begin{alertblock}{Đang phát triển}
        \footnotesize Tự động cho người chơi rời phòng khi mất mạng giữa chừng.
      \end{alertblock}
    \end{column}
  \end{columns}
\end{frame}

% ============================================================
% PHẦN: CƠ SỞ DỮ LIỆU VÀ GIAO DIỆN  (4 slide)
% Cấu trúc và ảnh chụp lấy từ nhánh dung/slides-phan-dung (Thành viên B: Dũng).
% Nội dung chữ đã ĐỐI CHIẾU với code thật và sửa các chỗ không đúng sự thật
% (tên bảng, cột Elo, "real-time", bên gọi REST...). Ảnh đều là ảnh chụp thật.
% ============================================================
\section{Cơ Sở Dữ Liệu Và Giao Diện}

% ---------- Slide: ERD ----------
\begin{frame}{Mô Hình Dữ Liệu (ERD)}
  \centering
  \begin{tikzpicture}[
    ent/.style={draw=inkblue, fill=white, rounded corners=2pt, thick, align=left,
                font=\scriptsize, inner sep=5pt},
    arr/.style={thick, color=brandblue}]
    \node[ent] (u) at (0,0)    {\textbf{Users}\\[2pt]\underline{Id} (PK)\\Username (UNIQUE)\\Password (BCrypt)\\Wins, Losses, Score\\Condition\\UserId (FK, bạn bè)};
    \node[ent] (m) at (5.4,0)  {\textbf{Matches}\\[2pt]\underline{Id} (PK)\\Player1Id (FK)\\Player2Id (FK)\\WinnerId (FK)\\StartTime, EndTime};
    \node[ent] (d) at (10.8,0) {\textbf{MatchDetails}\\[2pt]\underline{Id} (PK)\\MatchId (FK)\\PlayerId (FK)\\X\_Coordinate, Y\_Coordinate\\IsHit, TurnTime};

    \draw[arr] (u) -- node[above,font=\scriptsize]{1 -- N} (m);
    \draw[arr] (m) -- node[above,font=\scriptsize]{1 -- N} (d);

    \coordinate (ua) at ($(u.south)+(0,-1.1)$);
    \coordinate (da) at ($(d.south)+(0,-1.1)$);
    \draw[arr] (u.south) -- (ua) -- node[below,font=\scriptsize]{1 -- N} (da) -- (d.south);
  \end{tikzpicture}

  \vspace{0.3cm}
  \footnotesize
  Một người chơi có nhiều trận; một trận có nhiều lượt bắn (\textbf{MatchDetails}). Mật khẩu chỉ lưu dạng \textbf{BCrypt}.\\
  \textbf{Đang phát triển:} ghi kết quả trận và lịch sử lượt bắn vào hai bảng \textit{Matches}, \textit{MatchDetails}.
\end{frame}

% ---------- Slide: đăng nhập + chống dò mật khẩu ----------
\begin{frame}{Đăng Nhập Và Chống Dò Mật Khẩu}
  \begin{columns}[T]
    \begin{column}{0.48\textwidth}
      \centering
      \slideimg[0.85\textwidth]{04-client-login.png}{Màn hình đăng nhập}
      \par\vspace{2pt}\scriptsize Client gửi \texttt{LoginRequestDto} qua TCP; server đối chiếu với CSDL.
    \end{column}
    \begin{column}{0.48\textwidth}
      \centering
      \slideimg[0.85\textwidth]{07-login-lockout.png}{Thông báo khóa}
      \par\vspace{2pt}\scriptsize Sai 5 lần: server khóa 60 giây, Client hiện đếm ngược.
    \end{column}
  \end{columns}
  \vspace{0.25cm}
  \begin{exampleblock}{Ý nghĩa}
    \footnotesize Khóa theo IP và theo tài khoản: chặn dò mật khẩu hàng loạt và giảm truy vấn rác xuống CSDL.
  \end{exampleblock}
\end{frame}

% ---------- Slide: lobby ----------
\begin{frame}{Sảnh Chờ (Lobby)}
  \begin{columns}[T]
    \begin{column}{0.50\textwidth}
      \centering
      \slideimg[\textwidth]{05-client-lobby.png}{Giao diện Lobby}
    \end{column}
    \begin{column}{0.46\textwidth}
      \begin{block}{Đã hoạt động}
        \begin{itemize}
          \item \textbf{Join Room}: xem danh sách phòng từ server.
          \item \textbf{Leader Board}: bảng xếp hạng từ Supabase.
          \item \textbf{Profile}: thông tin người chơi.
        \end{itemize}
      \end{block}
      \begin{alertblock}{Đang phát triển}
        \footnotesize \textbf{Find Opponent} (ghép trận) và \textbf{Create Room}: đang nối vào giao diện.
      \end{alertblock}
    \end{column}
  \end{columns}
\end{frame}

% ---------- Slide: bảng xếp hạng ----------
\begin{frame}{Bảng Xếp Hạng}
  \begin{columns}[T]
    \begin{column}{0.52\textwidth}
      \begin{block}{Lấy dữ liệu qua web service}
        \begin{itemize}
          \item Client gọi \textbf{REST của Supabase} bằng \texttt{HttpClient}, đọc view \texttt{leaderboard}.
          \item Sắp xếp \textbf{điểm giảm dần}; bằng điểm thì ai thắng nhiều hơn đứng trước.
          \item Nút \textbf{Làm mới} tải lại; mất mạng thì báo lỗi nhẹ, \textbf{không làm sập ứng dụng}.
          \item View chỉ có 4 cột, \textbf{không có mật khẩu}.
        \end{itemize}
      \end{block}
    \end{column}
    \begin{column}{0.44\textwidth}
      \centering
      \slideimg[\textwidth]{06-client-leaderboard.png}{Form Bảng xếp hạng}
    \end{column}
  \end{columns}
\end{frame}

% PHẦN CỦA HOÀNG: Cloud, web service, đăng nhập an toàn. (Đức và Dũng KHÔNG sửa file này)
% =====================================================
\section{Cloud và Web Service}

\begin{frame}{Vì Sao Chọn Supabase?}
  \small
  \renewcommand{\arraystretch}{1.25}
  \begin{tabularx}{\textwidth}{@{}l Y Y Y@{}}
    \toprule
    \textbf{Tiêu chí} & \textbf{SQLite} & \textbf{Firebase} & \textbf{Supabase} \\
    \midrule
    Loại dữ liệu & File cục bộ (SQL) & NoSQL (Firestore) & \textbf{PostgreSQL} (SQL) \\
    Dùng được EF Core hiện có & Có & \textbf{Không} (không có provider) & \textbf{Có} (Npgsql) \\
    Web service có sẵn & Không & Có (SDK/REST) & \textbf{Có, REST tự sinh} \\
    Cần Internet & Không & Có & Có \\
    Công sức chuyển đổi & --- & Rất lớn (viết lại lớp dữ liệu) & \textbf{Nhỏ} (đổi provider) \\
    \bottomrule
  \end{tabularx}
  \vspace{0.3cm}
  \begin{exampleblock}{Kết luận}
    Supabase giữ nguyên mô hình quan hệ (Users, Matches, MatchDetails) và EF Core đã có, đồng thời cho \textbf{REST API} miễn phí. \textbf{SQLite vẫn giữ làm dự phòng} khi mất mạng.
  \end{exampleblock}
\end{frame}

\begin{frame}{Kiến Trúc Sau Khi Thêm Cloud}
  \centering
  \begin{tikzpicture}[
    box/.style={draw=inkblue, fill=white, rounded corners=3pt, thick, align=center,
                minimum height=1.1cm, minimum width=3.0cm, font=\footnotesize},
    cloud/.style={box, fill=paleblue},
    arr/.style={-{Stealth[length=2.2mm]}, thick, color=brandblue},
    lab/.style={font=\scriptsize, color=footernavy, fill=bgblue, inner sep=1pt}]
    \node[box]   (cli) at (0,0)    {\textbf{Client WinForms}\\Login, Lobby, Bàn cờ};
    \node[box]   (srv) at (5.2,0)  {\textbf{TCP Server C\#}\\Router, Handler};
    \node[cloud] (pg)  at (10.4,0) {\textbf{Supabase}\\PostgreSQL};
    \node[cloud] (api) at (10.4,-2.4) {\textbf{REST API}\\view \texttt{leaderboard}};
    \node[box]   (lite) at (5.2,2.3) {\textbf{SQLite}\\(dự phòng offline)};

    \draw[arr] (cli) -- node[lab,above] {TCP + AES-128} (srv);
    \draw[arr] (srv) -- node[lab,above] {EF Core + Npgsql} (pg);
    \draw[arr,dashed] (srv) -- node[lab,right] {đổi cấu hình} (lite);
    \draw[arr] (api) -- (pg);
    \draw[arr] (cli.south) |- node[lab,above,pos=0.75] {HTTPS + apikey} (api.west);
  \end{tikzpicture}
  \vspace{0.2cm}

  \footnotesize Đường game đi qua \textbf{TCP Server}; đường web service đi thẳng từ Client tới \textbf{REST của Supabase}.
\end{frame}

\begin{frame}{Kết Nối Cơ Sở Dữ Liệu Với EF Core}
  \begin{block}{Đổi nơi lưu trữ chỉ bằng cấu hình, không sửa code khi demo}
    \begin{itemize}
      \item \texttt{BATTLESHIP\_DB=supabase} $\rightarrow$ Server dùng \textbf{PostgreSQL} (Npgsql).
      \item \texttt{BATTLESHIP\_DB=sqlite} (mặc định) $\rightarrow$ dùng \texttt{battleship.db} cục bộ.
      \item Chuỗi kết nối đọc từ biến môi trường, \textbf{không nằm trong mã nguồn}.
      \item Server \textbf{tự tạo bảng} (\texttt{EnsureCreated}) khi khởi động.
    \end{itemize}
  \end{block}
  \begin{alertblock}{Lưu ý kỹ thuật}
    Dùng chuỗi kết nối \textbf{Session pooler} của Supabase (hỗ trợ IPv4) và \texttt{SSL Mode=Require}.
  \end{alertblock}
\end{frame}

\begin{frame}{Web Service: REST API Của Supabase}
  \begin{block}{Mỗi bảng/view trong Supabase tự có một endpoint}
    \texttt{GET https://<project>.supabase.co/rest/v1/leaderboard}\\
    \texttt{\ \ \ \ \ ?select=*\&order=Score.desc\&limit=10}
  \end{block}
  \begin{columns}[T]
    \begin{column}{0.48\textwidth}
      \begin{block}{Header bắt buộc}
        \begin{itemize}
          \item \texttt{apikey: <publishable key>} (bắt buộc)
          \item Thiếu khóa $\rightarrow$ \textbf{401}; khóa đúng nhưng chưa có bảng $\rightarrow$ 404
        \end{itemize}
      \end{block}
    \end{column}
    \begin{column}{0.48\textwidth}
      \begin{exampleblock}{Phản hồi (JSON)}
        \footnotesize
        \texttt{[\{"Username":"admiral\_a",}\\
        \texttt{\ \ "Wins":12,"Losses":3,}\\
        \texttt{\ \ "Score":1200\}, ...]}
      \end{exampleblock}
    \end{column}
  \end{columns}
  \vspace{0.1cm}
  \footnotesize Client C\# gọi bằng \texttt{HttpClient} và hiển thị lên \texttt{ListView} \textbf{Bảng xếp hạng}.
\end{frame}

\begin{frame}{Cơ Chế Đăng Nhập An Toàn}
  \footnotesize
  \renewcommand{\arraystretch}{1.2}
  \begin{tabularx}{\textwidth}{@{}Y Y@{}}
    \toprule
    \textbf{Nguy cơ} & \textbf{Cách phòng thủ} \\
    \midrule
    Gửi gói tin khi chưa đăng nhập & Router \textbf{chặn} mọi gói, trừ Handshake/Login/Register \\
    Dò mật khẩu hàng loạt & \textbf{Khóa 60 giây} sau 5 lần sai (theo IP và theo tài khoản) \\
    Dữ liệu đầu vào xấu & Kiểm tra tên 3--20 ký tự, mật khẩu 6--64 ký tự \\
    Một tài khoản đăng nhập nhiều nơi & Mỗi tài khoản \textbf{một phiên} tại một thời điểm \\
    Giả mạo \textit{Admin} và \textit{admin} & So sánh tên \textbf{không phân biệt hoa thường} \\
    Đoán tên tài khoản qua thời gian trả lời & Luôn chạy BCrypt, thông báo lỗi \textbf{chung chung} \\
    Lộ mật khẩu khi lưu trữ & Chỉ lưu chuỗi băm \textbf{BCrypt} (có muối) \\
    \bottomrule
  \end{tabularx}
  \vspace{0.15cm}
  \begin{exampleblock}{Kiểm thử}
    \footnotesize 19 kịch bản tấn công qua TCP, chạy trên server thật: \textbf{đều đạt}.
  \end{exampleblock}
\end{frame}

\begin{frame}{Bảo Mật Dữ Liệu}
  \begin{columns}[T]
    \begin{column}{0.56\textwidth}
      \begin{block}{Nhiều lớp bảo vệ}
        \footnotesize
        \begin{itemize}
          \item \textbf{Đường truyền game:} gói tin TCP mã hóa \textbf{AES-128}.
          \item \textbf{Mật khẩu:} chỉ lưu chuỗi băm \textbf{BCrypt}.
          \item \textbf{Supabase:} bật \textbf{Row Level Security}, thu quyền công khai trên bảng gốc.
          \item \textbf{Web service:} chỉ mở \textbf{view} \texttt{leaderboard}, không có cột mật khẩu.
          \item \textbf{Khóa:} Client chỉ giữ \texttt{publishable key}; mật khẩu DB nằm trong biến môi trường, \textbf{không commit}.
        \end{itemize}
      \end{block}
    \end{column}
    \begin{column}{0.40\textwidth}
      \centering
      \slideimg[\textwidth]{02-rest-blocked.png}{Nhật ký Supabase}
      \par\vspace{2pt}\scriptsize Đọc thẳng bảng \texttt{Users} bằng khóa công khai bị Postgres từ chối: \texttt{permission denied} (mã 42501).
    \end{column}
  \end{columns}
  \vspace{0.1cm}
  \begin{alertblock}{Vì sao phải khóa bảng gốc?}
    \footnotesize Supabase mặc định mở REST cho mọi bảng. Nếu không bật RLS, ai có \texttt{publishable key} đều đọc được cột mật khẩu băm của mọi người chơi.
  \end{alertblock}
\end{frame}

% PHẦN CỦA HOÀNG: Kịch bản demo và dự phòng. (Đức và Dũng KHÔNG sửa file này)
% =====================================================
\section{Kịch Bản Demo}

\begin{frame}{Demo Chứng Minh \textquotedblleft Móc Được Dữ Liệu\textquotedblright}
  \begin{enumerate}
    \item \textbf{Ghi:} đăng ký tài khoản trong game $\rightarrow$ mở \textbf{Supabase Table Editor}, thấy dòng mới, cột mật khẩu là chuỗi \texttt{\$2a\$...} (BCrypt).
    \item \textbf{Đọc qua web service:} gọi \texttt{GET .../rest/v1/leaderboard} $\rightarrow$ nhận \textbf{JSON}.
    \item \textbf{Đọc trong ứng dụng:} Lobby $\rightarrow$ nút \textbf{Bảng xếp hạng} hiển thị dữ liệu lấy qua REST.
    \item \textbf{Dự phòng:} đặt \texttt{BATTLESHIP\_DB=sqlite} $\rightarrow$ game vẫn chạy khi mất mạng.
  \end{enumerate}
  \vspace{0.2cm}
  \begin{exampleblock}{Thời lượng}
    Khoảng \textbf{5 phút}. Hoàng thuyết trình và điều khiển demo; Đức và Dũng hỗ trợ thao tác và trả lời câu hỏi thuộc mảng của mình.
  \end{exampleblock}
\end{frame}

\begin{frame}{Rủi Ro Và Phương Án Dự Phòng}
  \small
  \renewcommand{\arraystretch}{1.25}
  \begin{tabularx}{\textwidth}{@{}Y Y@{}}
    \toprule
    \textbf{Rủi ro} & \textbf{Phương án} \\
    \midrule
    Mạng trường chập chờn & Hotspot 4G; chuyển sang \textbf{SQLite}; phát \textbf{video demo} quay sẵn \\
    Supabase gói miễn phí tạm dừng project & Kiểm tra dashboard trước buổi demo, bấm \textbf{Restore} nếu cần \\
    Chuỗi kết nối \textit{direct} không nối được (IPv6) & Dùng chuỗi \textbf{Session pooler} \\
    Lộ khóa hoặc mật khẩu & Biến môi trường, không commit, không dùng \texttt{service\_role} ở Client \\
    \bottomrule
  \end{tabularx}
\end{frame}

% PHẦN CỦA HOÀNG: Phân công, kết luận, cảm ơn. (Đức và Dũng KHÔNG sửa file này)
% =====================================================
\section{Nhóm Thực Hiện}

\begin{frame}{Phân Công Thuyết Trình Và Thực Hiện}
  \small
  \renewcommand{\arraystretch}{1.3}
  \begin{tabularx}{\textwidth}{@{}>{\raggedright\arraybackslash}p{2.4cm} Y Y@{}}
    \toprule
    \textbf{Thành viên} & \textbf{Phần thuyết trình} & \textbf{Phần kỹ thuật} \\
    \midrule
    \textbf{Hoàng} \newline {\scriptsize \color{brandblue}\bevnsemi CORE THUYẾT TRÌNH}
      & Mở đầu, kiến trúc, bảo mật, \textbf{demo live}, kết luận
      & Supabase, chuyển DB (SQLite $\leftrightarrow$ Supabase), TCP/AES, Bảng xếp hạng qua REST \\
    \textbf{Dũng}
      & Mục \textit{Cơ sở dữ liệu} (ERD, RLS), thao tác máy demo
      & Models/EF Core, SQL \texttt{setup}, dữ liệu mẫu, nút Bảng xếp hạng ở Lobby \\
    \textbf{Đức}
      & Mục \textit{Kiến trúc server và game}, hỗ trợ thao tác demo
      & RoomManager, State Machine, video demo dự phòng \\
    \bottomrule
  \end{tabularx}
  \vspace{0.25cm}

  \footnotesize Hoàng nói chính; Đức và Dũng trả lời câu hỏi thuộc mảng của mình.
\end{frame}

\begin{frame}{Kết Luận}
  \begin{block}{Những gì đạt được}
    \begin{itemize}
      \item Game dùng \textbf{raw socket TCP} với giao thức tự thiết kế, mã hóa \textbf{AES-128}.
      \item \textbf{Đăng nhập an toàn:} BCrypt, chặn gói khi chưa đăng nhập, khóa khi dò mật khẩu.
      \item Dữ liệu lưu được trên \textbf{Supabase (PostgreSQL)} qua EF Core.
      \item Có \textbf{web service REST} để truy vấn bảng xếp hạng từ bên ngoài.
      \item Vẫn chạy được \textbf{offline} nhờ SQLite dự phòng.
    \end{itemize}
  \end{block}
  \begin{exampleblock}{Hướng phát triển}
    Đăng nhập bằng token của Supabase Auth, đồng bộ lịch sử trận đấu lên cloud, thêm bảng xếp hạng theo tuần.
  \end{exampleblock}
\end{frame}

\begin{frame}
  \centering
  \vspace{1.5cm}
  {\bungee\color{inkblue}\fontsize{20}{24}\selectfont CẢM ƠN ĐÃ THEO DÕI!}\par
  \vspace{0.3cm}
  {\bevnsemi\color{brandblue}\small Nhóm sẵn sàng trả lời câu hỏi}
\end{frame}


\end{document}
